\section{Phase 2 : Protocole de benchmark}
\label{sec:phase2}

\subsection{Jeux de caractéristiques}

À partir des paramètres optimisés en Phase~1, nous définissons huit jeux de
caractéristiques (Tableau~\ref{tab:feature_sets}). Les jeux mono-famille (A--D)
évaluent chaque descripteur isolément ; les jeux combinés testent si la fusion
ajoute une valeur discriminante.

\begin{table}[H]
  \centering
  \caption{Jeux de caractéristiques évalués en Phase~2.}
  \label{tab:feature_sets}
  \begin{tabular}{@{}lp{0.48\linewidth}c@{}}
    \toprule
    \textbf{Label} & \textbf{Contenu} & \textbf{Dim. nettoyée} \\
    \midrule
    A --- Stats       & Statistiques du premier ordre uniquement   & 9 \\
    B --- Tex(opt)    & Texture LBP optimisée                      & 18 \\
    C --- DWT(opt)    & Statistiques des sous-bandes DWT optimisées & 12 \\
    D --- HOG         & Histogrammes HOG forme/contours            & 490 \\
    A+B               & Statistiques + LBP                         & 27 \\
    B+C(opt)          & LBP + DWT                                  & 30 \\
    A+B+C(opt)        & Statistiques + LBP + DWT                   & 36 \\
    Full(opt)         & Les quatre familles (A+B+C+D)              & 526 \\
    \bottomrule
  \end{tabular}
\end{table}

Chaque jeu de caractéristiques subit un nettoyage post-concaténation : un filtre
de variance supprime les caractéristiques quasi nulles, suivi d'un élagage par
corrélation (seuil $r > 0{,}95$), puis d'une normalisation StandardScaler dans
chaque \texttt{Pipeline} sklearn.

\subsection{Classifieurs}

Huit classifieurs ont été benchmarkés :

\begin{table}[H]
  \centering
  \caption{Classifieurs et leur rôle dans le benchmark.}
  \label{tab:classifiers}
  \begin{tabular}{@{}lll@{}}
    \toprule
    \textbf{Famille} & \textbf{Classifieur} & \textbf{Propriété clé} \\
    \midrule
    Marge / noyau & SVM (RBF / polynomial) & frontières non linéaires \\
    Basé sur instances & KNN & distances locales, sans entraînement \\
    Arbres bagués & Random Forest, ExtraTrees & réduction de variance par ensemble \\
    Arbres boostés & XGBoost, LightGBM & correction séquentielle des erreurs \\
    Linéaire & Régression logistique & référence rapide et interprétable \\
    Réseau peu profond & MLP (deux couches cachées) & activation non linéaire \\
    \bottomrule
  \end{tabular}
\end{table}

\subsection{Protocole d'évaluation}

Chaque classifieur est ajusté par \textbf{GridSearchCV à 5 plis} sur l'ensemble
d'entraînement, puis évalué sur la partition de test tenue à l'écart. Six modes
d'évaluation ont été suivis ; les résultats principaux rapportés ici correspondent
à la partition de test. Cela donne
\NModels{}~classifieurs $\times$ \NFeatureSets{}~jeux de caractéristiques = \textbf{\NRuns{}~pipelines}.

Métrique principale : F1-macro. Métriques secondaires : précision (accuracy),
précision et rappel macro-moyennés, et $\kappa$ de Cohen.
